\documentclass[11pt,a4paper]{article}

% ── Packages ─────────────────────────────────────────────────────────────────
\usepackage[utf8]{inputenc}
\usepackage[T1]{fontenc}
\usepackage{amsmath,amssymb,amsthm}
\usepackage{mathtools}
\usepackage{physics}
\usepackage{hyperref}
\usepackage{cleveref}
\usepackage{microtype}
\usepackage{geometry}
\geometry{margin=2.5cm}
\usepackage{booktabs}

% ── Theorem environments ──────────────────────────────────────────────────────
\newtheorem{theorem}{Theorem}[section]
\newtheorem{lemma}[theorem]{Lemma}
\newtheorem{corollary}[theorem]{Corollary}
\newtheorem{definition}[theorem]{Definition}
\newtheorem{remark}[theorem]{Remark}

% ── Macros ───────────────────────────────────────────────────────────────────
\newcommand{\F}{\mathbb{F}}
\newcommand{\Z}{\mathbb{Z}}
\newcommand{\rk}{\operatorname{rank}_{\F_2}}
\newcommand{\VNE}{S}
\newcommand{\supp}{\operatorname{supp}}

% ─────────────────────────────────────────────────────────────────────────────
\title{Analytic von Neumann Entropy for CNOT Circuits\\
via Walsh–Hadamard Transforms over $\F_2$}

\author{\"Ozg\"ur \"Ozer\\
\small \texttt{[email]}}

\date{\today}

\begin{document}
\maketitle

% ─────────────────────────────────────────────────────────────────────────────
\begin{abstract}
We derive a closed-form expression for the bipartite von Neumann entropy of
pure states produced by CNOT circuits acting on product inputs.
For a circuit in which $n_A$ qubits, each prepared in $R_x(\theta_i)|0\rangle$,
are coupled to $n_B$ ancilla qubits via a CNOT layer encoded by a binary
matrix $M \in \F_2^{n_B \times n_A}$, the entropy of the ancilla subsystem
equals the Shannon entropy of an explicit probability distribution
obtained by a Walsh–Hadamard transform (Theorem~\ref{thm:wh}).
This gives tight algebraic bounds: $\VNE(B) \leq \rk(M)$, with equality at
$\theta_i = \pi/2$. We extend the formula to multi-layer circuits, arbitrary
pure inputs, and diagonal mixed states, and prove that a bidirectional CNOT
circuit satisfies $U^3 = I$ over $\F_2^2$, yielding period-3 entropy dynamics.
All results are accompanied by an open-source numerical verification library
available at \url{https://github.com/Ozgurooozer/qse}.
\end{abstract}

\tableofcontents
\bigskip

% ─────────────────────────────────────────────────────────────────────────────
\section{Introduction}
\label{sec:intro}

Understanding the structure of entanglement in quantum circuits is a central
problem in quantum information theory.  For stabilizer states generated by
Clifford circuits, bipartite entanglement is known to be determined by
$\F_2$-linear algebraic invariants, such as the rank of associated binary
matrices~\cite{audenaert2005entanglement,fattal2004entanglement}.  However,
these results apply primarily to discrete families of states, where
entanglement takes quantized values.  In contrast, many physically relevant
quantum circuits involve continuous parameters, such as rotation angles, for
which analytic expressions for the von Neumann entropy are generally not
available.

In this work, we study a restricted but natural class of parameterized quantum
circuits: product states of the form $\bigotimes_i R_x(\theta_i)\ket{0}$ on
subsystem $A$, coupled via a layer of CNOT gates to an initially unentangled
subsystem $B$.  While the entangling layer is Clifford, the input state is
not, except at special angles, placing the resulting states outside the
stabilizer formalism.

Our main contribution is an explicit analytic expression for the bipartite von
Neumann entropy of subsystem $B$ as a function of the rotation angles
$\{\theta_i\}$ and the binary connectivity matrix $M$.  Specifically, we show
that the entropy can be written as the Shannon entropy of a probability
distribution obtained via a Walsh--Hadamard transform over $\F_2^{n_B}$
(Theorem~\ref{thm:wh}).

This result recovers the known rank-based entropy formula in the stabilizer
limit $\theta_i = \pi/2$, while extending it to a continuous family of
non-stabilizer states.  We further show that several structural properties of
the circuit---including entropy bounds, optimality conditions, and multi-layer
composition---admit simple descriptions in terms of $\F_2$ linear algebra.

\paragraph{Scope and limitations.}
Our analysis applies to circuits consisting of unidirectional CNOT layers
($A \to B$) acting on product inputs, with subsystem $B$ initialized in
$\ket{0}^{\otimes n_B}$.  Circuits involving bidirectional interactions,
additional Clifford gates, or entangled inputs on $B$ fall outside this
framework.

\paragraph{Contributions.}
The main contribution of this work is Theorem~\ref{thm:wh}, which provides a
closed-form Walsh--Hadamard expression for the bipartite von Neumann entropy in
a class of parameterized CNOT circuits.  Secondary contributions include
rank-based entropy bounds, optimality conditions, and structural results on
circuit composition.

\paragraph{Relation to prior work.}
Our results build on the stabilizer formalism, where entanglement is
characterized by algebraic invariants of binary
matrices~\cite{audenaert2005entanglement,fattal2004entanglement}.  Unlike
these discrete results, our formula provides a closed-form expression for the
continuous dependence of entropy on rotation parameters.  The use of
Walsh--Hadamard transforms is closely related to Fourier analysis over
$\F_2^n$, which has appeared in studies of stabilizer entropy and related
quantities~\cite{bataille2022cnot,goodenough2024stabilizer}.

\subsection*{Related work}

The connection between stabiliser circuits and $\F_2$ linear algebra is
classical~\cite{gottesman1997stabilizer}.  Fattal et
al.~\cite{fattal2004entanglement} showed that the entanglement entropy of a
stabilizer state across any bipartition equals the $\F_2$-rank of a submatrix
of the stabilizer generators.  Audenaert and
Plenio~\cite{audenaert2005entanglement} proved that every bipartite mixed
stabilizer state is locally equivalent to a direct product of maximally
entangled pairs and a separable state.  These results establish the rank
formula $\VNE = \rk(M)$ at $\theta = \pi/2$, which our optimality theorem
recovers as a special case.  Bataille~\cite{bataille2022cnot} studied
entanglement creation by CNOT circuits acting on factorized states; our
multi-layer composition rule (Theorem~\ref{thm:layers}) is consistent with
that picture.  Goodenough et al.~\cite{goodenough2024stabilizer} connected
entanglement entropy of noisy stabilizer states to syndrome entropy, which is
complementary to our continuous-angle approach.  Entanglement in parameterised
circuits has also been studied in the context of
expressibility~\cite{sim2019expressibility} and barren
plateaus~\cite{mcclean2018barren}; the Walsh--Hadamard approach here gives
analytic control absent in those statistical treatments.

% ─────────────────────────────────────────────────────────────────────────────
\section{Setup}
\label{sec:prelim}

\paragraph{Notation.}
$\F_2 = \{0,1\}$ is the field with two elements; arithmetic is mod~2.
For $x \in \{0,1\}^n$, $|x\rangle$ denotes the corresponding
computational-basis state.  We write $\rk(M)$ for the rank of
$M \in \F_2^{n_B \times n_A}$ over $\F_2$.

\paragraph{Circuit.}
Let $n_A, n_B \geq 1$.  We prepare the input state
\begin{equation}
  |\psi_{\mathrm{in}}\rangle
  = \Bigl(\bigotimes_{i=0}^{n_A-1} R_x(\theta_i)|0\rangle\Bigr)
    \otimes |0\rangle^{\otimes n_B},
  \label{eq:input}
\end{equation}
where $R_x(\theta)|0\rangle = \cos(\theta/2)|0\rangle - i\sin(\theta/2)|1\rangle$.
A layer of CNOT gates is then applied: for each $(i,j)$ with $M_{j,i}=1$,
gate $\mathrm{CNOT}(A_i \to B_j)$ is applied.  We assume that the CNOT
gates within a layer act on disjoint target qubits (at most one $1$ per row
of $M$), so their order is immaterial; this assumption is relaxed in
Section~\ref{sec:layers}.

The output state is $|\psi\rangle = U_M|\psi_{\mathrm{in}}\rangle$ and we study
\begin{equation}
  \VNE(B) = -\operatorname{tr}[\rho_B \log_2 \rho_B],
  \qquad \rho_B = \operatorname{tr}_A[|\psi\rangle\langle\psi|].
\end{equation}

% ─────────────────────────────────────────────────────────────────────────────
\section{The Walsh–Hadamard Entropy Formula}
\label{sec:wh}

\begin{theorem}[Walsh–Hadamard formula]
\label{thm:wh}
Let $|\psi\rangle$ be the output of the circuit in Section~\ref{sec:prelim}.
For each $b \in \{0,1\}^{n_B}$, define
\begin{equation}
  p_b = \frac{1}{2^{n_B}}
        \sum_{s \in \{0,1\}^{n_B}}
        (-1)^{b \cdot s}\, \varphi(s),
  \label{eq:wh}
\end{equation}
where $b \cdot s = \sum_j b_j s_j \pmod{2}$ and
\begin{equation}
  \varphi(s) = \prod_{\substack{i=0\\(M^\top s)_i = 1}}^{n_A-1} \cos\theta_i.
  \label{eq:phi}
\end{equation}
(An empty product equals $1$.)
Then $p_b \geq 0$ for all $b$, $\sum_b p_b = 1$, and
\begin{equation}
  \boxed{\VNE(B) = H\!\bigl(\{p_b\}_{b \in \{0,1\}^{n_B}}\bigr),}
  \label{eq:entropy}
\end{equation}
where $H(\{q_b\}) = -\sum_b q_b \log_2 q_b$ is the Shannon entropy.
\end{theorem}

\begin{proof}
\textbf{Step 1: Structure of $\rho_B$.}
Expanding $R_x(\theta_i)|0\rangle = \cos(\theta_i/2)|0\rangle - i\sin(\theta_i/2)|1\rangle$,
the input state decomposes as
\begin{equation}
  |\psi_{\mathrm{in}}\rangle
  = \sum_{x \in \{0,1\}^{n_A}} \alpha_x\, |x\rangle_A\, |0\rangle_B,
  \quad
  \alpha_x = \prod_{i=0}^{n_A-1}
  \bigl(\cos\tfrac{\theta_i}{2}\bigr)^{1-x_i}
  \bigl(-i\sin\tfrac{\theta_i}{2}\bigr)^{x_i}.
  \label{eq:alpha}
\end{equation}
The CNOT layer maps $|x\rangle_A|0\rangle_B \mapsto |x\rangle_A|Mx\rangle_B$
(with arithmetic in $\F_2$), giving
\begin{equation}
  |\psi\rangle = \sum_{x \in \{0,1\}^{n_A}} \alpha_x\, |x\rangle_A\, |Mx\rangle_B.
  \label{eq:output}
\end{equation}
The reduced density matrix is
\begin{equation}
  \rho_B
  = \operatorname{tr}_A\bigl[|\psi\rangle\langle\psi|\bigr]
  = \sum_{b,b' \in \{0,1\}^{n_B}}
    \Bigl(\sum_{x:\, Mx=b} \alpha_x\Bigr)
    \Bigl(\sum_{x':\, Mx'=b'} \alpha_{x'}\Bigr)^*
    |b\rangle\langle b'|.
\end{equation}
For $b \neq b'$, the cross terms $\langle b|b'\rangle = 0$, so
$\rho_B$ is diagonal:
\begin{equation}
  \rho_B = \sum_{b \in \{0,1\}^{n_B}} \lambda_b\, |b\rangle\langle b|,
  \qquad
  \lambda_b = \sum_{x:\, Mx = b} |\alpha_x|^2 \geq 0.
  \label{eq:eigenvalues}
\end{equation}
Note that $|\alpha_x|^2 = \prod_i \cos^2(\theta_i/2)^{1-x_i} \sin^2(\theta_i/2)^{x_i}$,
and $\sum_b \lambda_b = \sum_x |\alpha_x|^2 = 1$ by the product structure.

\textbf{Step 2: Walsh–Hadamard inversion.}
Regard $\lambda : \{0,1\}^{n_B} \to [0,1]$ as a function on $\F_2^{n_B}$.
Define its Walsh–Hadamard transform
$\hat\lambda(s) = \sum_b (-1)^{b\cdot s}\lambda_b$.
By Fourier inversion on $\F_2^{n_B}$,
\begin{equation}
  \lambda_b = \frac{1}{2^{n_B}} \sum_{s} (-1)^{b\cdot s} \hat\lambda(s).
\end{equation}
It remains to evaluate $\hat\lambda(s)$.  Using~\eqref{eq:eigenvalues},
\begin{align}
  \hat\lambda(s)
  &= \sum_b (-1)^{b\cdot s} \sum_{x:\,Mx=b} |\alpha_x|^2
   = \sum_x (-1)^{(Mx)\cdot s} |\alpha_x|^2 \notag\\
  &= \sum_x (-1)^{x\cdot M^\top s} |\alpha_x|^2
   = \prod_{i=0}^{n_A-1}
     \sum_{x_i \in \{0,1\}} (-1)^{x_i (M^\top s)_i}
     \cos^2(\theta_i/2)^{1-x_i} \sin^2(\theta_i/2)^{x_i}.
   \label{eq:hat_lambda}
\end{align}
The inner sum over each $x_i$ evaluates to
\begin{equation}
  \sum_{x_i \in \{0,1\}} (-1)^{x_i c_i} p_i^{1-x_i} q_i^{x_i}
  = p_i - (-1)^{c_i} q_i
  \quad (c_i = (M^\top s)_i,\; p_i = \cos^2(\theta_i/2),\; q_i = \sin^2(\theta_i/2)).
\end{equation}
Substituting $p_i = \cos^2(\theta_i/2)$, $q_i = \sin^2(\theta_i/2)$:
\begin{itemize}
  \item If $c_i = 0$: sum $= p_i + q_i = 1$.
  \item If $c_i = 1$: sum $= p_i - q_i = \cos^2(\theta_i/2) - \sin^2(\theta_i/2) = \cos\theta_i$.
\end{itemize}
Therefore $\hat\lambda(s) = \prod_{i:\,(M^\top s)_i=1} \cos\theta_i = \varphi(s)$,
and the Fourier inversion formula gives $\lambda_b = p_b$ as
in~\eqref{eq:wh}.  Since $\lambda_b \geq 0$ and $\sum_b \lambda_b = 1$,
the $p_b$ form a probability distribution and
$\VNE(B) = -\sum_b \lambda_b \log_2 \lambda_b = H(\{p_b\})$.
\end{proof}

\begin{remark}[Complexity]
Computing $\{\varphi(s)\}_{s \in \{0,1\}^{n_B}}$ takes $O(n_A \cdot 2^{n_B})$ time;
the Walsh--Hadamard inversion then takes $O(n_B \cdot 2^{n_B})$ time.
The naive implementation of~\eqref{eq:wh} iterates over all $b$ and $s$,
giving $O(4^{n_B} \cdot n_B)$ overall.
Both are exponential in $n_B$, but the formula avoids the full $2^{n_A+n_B}$
state-vector computation when $n_A \gg n_B$.  Empirical runtimes confirm the
$O(4^{n_B})$ scaling: the ratio $t(n_B{+}1)/t(n_B)$ converges to $\approx 4$
for $n_B \geq 5$ (see \texttt{tests/test\_scaling.py}).
\end{remark}

% ─────────────────────────────────────────────────────────────────────────────
\section{Rank Bounds}
\label{sec:rank}

\begin{theorem}[Entropy rank bound]
\label{thm:rank}
For all $\theta \in \mathbb{R}^{n_A}$ and $M \in \F_2^{n_B \times n_A}$,
\begin{equation}
  \VNE(B) \leq \rk(M).
  \label{eq:rank_bound}
\end{equation}
\end{theorem}

\begin{proof}
From~\eqref{eq:eigenvalues}, $\rho_B$ is supported on
$\mathrm{Im}_{\F_2}(M) = \{Mx : x \in \F_2^{n_A}\}$,
which has cardinality $2^{\rk(M)}$.
Hence $\rho_B$ has at most $2^{\rk(M)}$ nonzero eigenvalues, and
$\VNE(B) \leq \log_2(2^{\rk(M)}) = \rk(M)$.
\end{proof}

\begin{theorem}[Optimality at $\theta = \pi/2$]
\label{thm:opt}
If $\theta_i = \pi/2$ for all $i$, then $\VNE(B) = \rk(M)$.
\end{theorem}

\begin{proof}
At $\theta_i = \pi/2$, we have $|\alpha_x|^2 = 2^{-n_A}$ for all
$x \in \{0,1\}^{n_A}$, since
$\cos^2(\pi/4) = \sin^2(\pi/4) = 1/2$.
Each coset of $\ker_{\F_2}(M)$ has size $|\ker_{\F_2}(M)| = 2^{n_A - \rk(M)}$,
so from~\eqref{eq:eigenvalues},
\begin{equation}
  \lambda_b =
  \begin{cases}
    2^{-\rk(M)} & b \in \mathrm{Im}_{\F_2}(M), \\
    0            & \text{otherwise.}
  \end{cases}
\end{equation}
Hence $\VNE(B) = \log_2(2^{\rk(M)}) = \rk(M)$.
\end{proof}

% ─────────────────────────────────────────────────────────────────────────────
\section{Multi-Layer Circuits}
\label{sec:layers}

\begin{theorem}[$\F_2$-composition of layers]
\label{thm:layers}
Let $M_1, \ldots, M_k \in \F_2^{n_B \times n_A}$ be connectivity matrices of
$k$ sequential CNOT layers applied to the same input state~\eqref{eq:input}.
Then
\begin{equation}
  \VNE(B) = \VNE^{(\mathrm{Thm.\,\ref{thm:wh}})}\!\bigl(\theta,\; M_{\mathrm{eff}}\bigr),
  \qquad
  M_{\mathrm{eff}} = \sum_{\ell=1}^{k} M_\ell \pmod{2}.
  \label{eq:meff}
\end{equation}
\end{theorem}

\begin{proof}
Each $\mathrm{CNOT}(A_i \to B_j)$ gate is its own inverse.
For a fixed pair $(i,j)$, the gate appears $\sum_\ell (M_\ell)_{j,i}$ times;
an even count is equivalent to no gate, an odd count to one gate.
Hence the composition of all $k$ layers is unitarily equivalent to the
single layer with matrix $M_{\mathrm{eff}} = \sum_\ell M_\ell \pmod{2}$.
Applying Theorem~\ref{thm:wh} to $M_{\mathrm{eff}}$ gives~\eqref{eq:meff}.
\end{proof}

\begin{remark}
Theorem~\ref{thm:layers} does not require disjoint targets per layer;
it only requires that the \emph{composition} be well-defined, which holds
for any sequence of CNOT layers.
\end{remark}

\begin{corollary}[Period-2 sequence]
\label{cor:period2}
If all layers use the same matrix $M$, the entropy sequence satisfies
$\VNE^{(k)} = \VNE^{(k+2)}$ for all $k \geq 1$.
\end{corollary}

\begin{proof}
$M_{\mathrm{eff}}(k) = k \cdot M \pmod{2}$, which equals $M$ for odd $k$
and $0$ for even $k$.  When $M_{\mathrm{eff}} = 0$, the B-system remains in
$|0\rangle^{\otimes n_B}$ and $\VNE(B) = 0$.
\end{proof}

% ─────────────────────────────────────────────────────────────────────────────
\section{General Pure-State Input}
\label{sec:pure}

\begin{theorem}[Arbitrary pure input on $A$]
\label{thm:pure}
Let $|\psi_A\rangle = \sum_{x \in \{0,1\}^{n_A}} \alpha_x |x\rangle$ be any
normalised pure state on $A$, and let $B$ be initialised in
$|0\rangle^{\otimes n_B}$.  After applying $U_M$,
\begin{equation}
  \VNE(B) = H\!\Bigl(\bigl\{q_b\bigr\}_{b \in \{0,1\}^{n_B}}\Bigr),
  \qquad
  q_b = \sum_{x:\, Mx = b} |\alpha_x|^2.
  \label{eq:pure}
\end{equation}
\end{theorem}

\begin{proof}
Steps~1 and 2 of the proof of Theorem~\ref{thm:wh} apply verbatim with
$|\alpha_x|^2 = |\langle x|\psi_A\rangle|^2$ in place of the product form.
The diagonality of $\rho_B$ and the eigenvalue formula~\eqref{eq:eigenvalues}
hold for any product-basis decomposition of $|\psi_A\rangle$.
\end{proof}

\begin{corollary}[Blindness condition]
\label{cor:blind}
$\VNE(B) = 0$ if and only if $M(x \oplus x') = 0$ for all
$x, x' \in \supp(|\psi_A\rangle)$, where
$\supp(|\psi_A\rangle) = \{x : \alpha_x \neq 0\}$.
\end{corollary}

\begin{proof}
$\VNE(B) = 0$ if and only if all $q_b \in \{0, 1\}$, i.e.\ the map
$x \mapsto Mx$ is constant on $\supp(|\psi_A\rangle)$.  This is equivalent
to $M(x \oplus x') = Mx \oplus Mx' = 0$ for all pairs in the support.
\end{proof}

% ─────────────────────────────────────────────────────────────────────────────
\section{Diagonal Mixed Input}
\label{sec:mixed}

\begin{theorem}[Diagonal mixed input]
\label{thm:mixed}
Suppose $A$ is in a state $\rho_A$ that is diagonal in the computational
basis: $\rho_A = \sum_{x \in \{0,1\}^{n_A}} d_x |x\rangle\langle x|$,
$d_x \geq 0$, $\sum_x d_x = 1$.  After applying $U_M$,
\begin{equation}
  \VNE(B) = H\!\Bigl(\bigl\{r_b\bigr\}\Bigr),
  \qquad
  r_b = \sum_{x:\, Mx = b} d_x.
  \label{eq:mixed}
\end{equation}
\end{theorem}

\begin{proof}
Write $\rho_{AB} = \sum_x d_x\, U_M(|x\rangle\langle x| \otimes |0\rangle\langle 0|)U_M^\dagger$.
Since $U_M$ maps $|x\rangle_A|0\rangle_B$ to $|x\rangle_A|Mx\rangle_B$,
tracing out $A$ gives $\rho_B = \sum_b r_b |b\rangle\langle b|$
with $r_b$ as in~\eqref{eq:mixed}.
\end{proof}

\begin{remark}
The off-diagonal entries of $\rho_A$ play no role because $U_M$ is diagonal
in the $A$-basis: $\langle x|U_M|x'\rangle_A = 0$ for $x \neq x'$.
If $\rho_A$ is \emph{not} diagonal in the computational basis,
the off-diagonal coherences generally affect $\VNE(B)$; see Theorem~\ref{thm:pure}
for the pure-state case.
\end{remark}

% ─────────────────────────────────────────────────────────────────────────────
\section{CZ-Gate Equivalence}
\label{sec:cz}

\begin{theorem}[CZ equivalence]
\label{thm:cz}
Consider the circuit obtained by replacing each $\mathrm{CNOT}(A_i \to B_j)$
by a $\mathrm{CZ}(A_i, B_j)$ gate and initialising each $B$-qubit in
$|+\rangle = (|0\rangle + |1\rangle)/\sqrt{2}$ instead of $|0\rangle$.
The von Neumann entropy of $B$ is unchanged:
\begin{equation}
  \VNE^{\mathrm{CZ}}(B) = \VNE^{\mathrm{CNOT}}(B).
\end{equation}
\end{theorem}

\begin{proof}
For each target qubit $B_j$, the conjugation identity
$\mathrm{CZ} = (I_A \otimes H_{B_j})\,\mathrm{CNOT}(A \to B_j)\,(I_A \otimes H_{B_j})$
shows that the CZ circuit with $B_j$ initialised in $H|0\rangle = |+\rangle$
equals $(I \otimes H^{\otimes n_B})$ applied to the CNOT output.
Since $H^{\otimes n_B}$ is a local unitary on $B$, it does not change the
spectrum of $\rho_B$, and hence leaves $\VNE(B)$ invariant.
\end{proof}

% ─────────────────────────────────────────────────────────────────────────────
\section{Bidirectional CNOT and Period-3 Dynamics}
\label{sec:bidir}

We study a $2+2$ qubit circuit defined as follows.
Label the qubits $A_0, A_1, B_0, B_1$.
Define the \emph{forward layer}
$U_{\mathrm{fwd}} = \mathrm{CNOT}(A_0 \to B_0)\,\mathrm{CNOT}(A_1 \to B_1)$
and the \emph{backward layer}
$U_{\mathrm{bwd}} = \mathrm{CNOT}(B_0 \to A_0)\,\mathrm{CNOT}(B_1 \to A_1)$.
Set $U = U_{\mathrm{bwd}}\,U_{\mathrm{fwd}}$.

\begin{theorem}[Period-3 dynamics]
\label{thm:period3}
The unitary $U$ acts on computational-basis states as
\begin{equation}
  U\,|a_0, a_1, b_0, b_1\rangle
  = |b_0,\; b_1,\; a_0 \oplus b_0,\; a_1 \oplus b_1\rangle.
  \label{eq:U_action}
\end{equation}
Treating the pair $(a, b) \in \F_2^2$ as a single register (where
$a = (a_0, a_1)$ and $b = (b_0, b_1)$ act independently on each
coordinate), $U$ implements the map $(a, b) \mapsto (b,\, a \oplus b)$,
which corresponds to multiplication by the matrix
$T = \bigl(\begin{smallmatrix} 0 & 1 \\ 1 & 1 \end{smallmatrix}\bigr)$
over $\F_2$.
Since $T^3 = I$ over $\F_2$ (the order of $T$ in $\mathrm{GL}_2(\F_2)$ is $3$),
we have $U^3 = I_{16}$ but $U^2 \neq I_{16}$.
Consequently, $\VNE^{(k)}(B) = \VNE^{(k+3)}(B)$ for all $k \geq 1$,
with minimal period $3$.
\end{theorem}

\begin{proof}
\textbf{Action of $U$.}
$U_{\mathrm{fwd}}$ maps $(a_0, a_1, b_0, b_1) \mapsto (a_0, a_1, b_0 \oplus a_0, b_1 \oplus a_1)$.
Denoting the result $(a_0, a_1, b_0', b_1')$ with $b_j' = b_j \oplus a_j$,
$U_{\mathrm{bwd}}$ then maps
$(a_0, a_1, b_0', b_1') \mapsto (a_0 \oplus b_0', a_1 \oplus b_1', b_0', b_1')
= (b_0, b_1, b_0 \oplus a_0, b_1 \oplus a_1)$,
establishing~\eqref{eq:U_action}.

\textbf{Order of $U$.}
Treating each coordinate pair $(a_j, b_j) \in \F_2^2$ independently,
$U$ acts as left-multiplication by
$T = \bigl(\begin{smallmatrix} 0 & 1 \\ 1 & 1 \end{smallmatrix}\bigr) \in \mathrm{GL}_2(\F_2)$.
One verifies:
\begin{equation}
  T^2 = \begin{pmatrix} 1 & 1 \\ 1 & 0 \end{pmatrix}, \qquad
  T^3 = \begin{pmatrix} 1 & 0 \\ 0 & 1 \end{pmatrix} = I_2
  \pmod{2}.
\end{equation}
Since $T^2 \neq I_2$ and $T^3 = I_2$, the order of $T$ is $3$, so $U^3 = I_{16}$
and $U^2 \neq I_{16}$.

\textbf{Entropy period.}
Since $U^3 = I_{16}$, the state after $k$ full rounds satisfies
$|\psi^{(k)}\rangle = |\psi^{(k+3)}\rangle$, hence
$\rho_B^{(k)} = \rho_B^{(k+3)}$ and $\VNE^{(k)} = \VNE^{(k+3)}$.
The period is minimal because generically
$\VNE^{(1)} \neq \VNE^{(2)}$ (verified numerically for generic $\theta$).
\end{proof}

\begin{remark}
The matrix $T = \bigl(\begin{smallmatrix} 0 & 1 \\ 1 & 1 \end{smallmatrix}\bigr)$
is a generator of $\mathrm{GL}_2(\F_2) \cong S_3$; the order-3 subgroup it
generates corresponds to the even permutations $\{e, (123), (132)\}$.
The period-3 dynamics thus reflects the group structure of $\mathrm{GL}_2(\F_2)$,
rather than any special property of the entropy function.
\end{remark}

% ─────────────────────────────────────────────────────────────────────────────
\section{Numerical Implementation and Verification}
\label{sec:implementation}

All theorems are verified numerically by the \texttt{qse} Python library,
available at \url{https://github.com/Ozgurooozer/qse}.
The library depends only on NumPy and can be installed via
\texttt{pip install .}.  A standalone script reproduces all tests:
\begin{verbatim}
  git clone https://github.com/Ozgurooozer/qse
  cd qse && python quick_proofs.py
\end{verbatim}

Each test compares the analytic formula against an independent
exact state-vector simulation on randomly sampled parameters.
Results are summarised in Table~\ref{tab:results}.

\begin{table}[h]
\centering
\begin{tabular}{lllcc}
\toprule
Theorem & Statement & Samples & Result \\
\midrule
\ref{thm:wh}     & Walsh–Hadamard formula            & 50  & $\max\Delta < 10^{-13}$ \\
\ref{thm:rank}   & $\VNE(B) \leq \rk(M)$            & 200 & 0 violations            \\
\ref{thm:opt}    & $\VNE(\pi/2) = \rk(M)$           & 100 & $\max\Delta < 10^{-13}$ \\
\ref{thm:layers} & Multi-layer XOR composition       & 50  & $\max\Delta < 10^{-13}$ \\
\ref{thm:mixed}  & Diagonal mixed input              & 40  & $\max\Delta < 2\times10^{-13}$ \\
\ref{thm:pure}   & Arbitrary pure input              & 50  & $\max\Delta < 2\times10^{-13}$ \\
\ref{cor:blind}  & Blindness condition               & 50  & 50/50 exact             \\
\ref{cor:period2}& Period-2 sequence                 & 100 & $\max\Delta = 0$        \\
\ref{thm:period3}& $U^3 = I_{16}$                   & —   & $\max|U^3 - I| = 0$    \\
\ref{thm:period3}& Period-3 entropy                 & 20  & $\max\Delta = 0$        \\
\ref{thm:cz}     & CZ equivalence                    & 30  & $\max\Delta < 10^{-13}$ \\
\bottomrule
\end{tabular}
\caption{Numerical verification (random seed 42; pass threshold $10^{-10}$;
$\Delta$ denotes the absolute error between formula and state-vector simulation).}
\label{tab:results}
\end{table}

% ─────────────────────────────────────────────────────────────────────────────
\section{Discussion}
\label{sec:discussion}

We have shown that the bipartite von Neumann entropy of CNOT circuits with
product inputs is controlled entirely by the $\F_2$-linear structure of the
connectivity matrix $M$ and the cosines of the rotation angles.  Several
directions remain open.

\paragraph{Differentiability and variational algorithms.}
The formula in Theorem~\ref{thm:wh} is smooth in $\theta$, and its gradient
is computable in closed form.  This makes it a natural analytical testbed for
studying entropy in parameterised circuits, complementing existing statistical
approaches to barren plateaus~\cite{mcclean2018barren}.

\paragraph{Entangled inputs.}
Theorem~\ref{thm:pure} handles arbitrary pure inputs but does not give a
formula in terms of circuit parameters; extending the Walsh–Hadamard approach
to entangled inputs requires understanding how $\F_2$-linearity interacts with
off-diagonal coherences.

\paragraph{General Clifford circuits.}
Adding Hadamard and phase gates preserves the $\F_2$-linear structure of the
stabiliser group, but the diagonal form of $\rho_B$ may be lost.
A Walsh-spectrum formula for general Clifford entropy is an open problem.

\paragraph{Tensor networks.}
The period-3 result (Theorem~\ref{thm:period3}) shows that the bidirectional
CNOT circuit implements a map of order 3 in $\mathrm{GL}_2(\F_2)$.
Understanding this in the tensor-network language may yield efficient classical
simulation schemes beyond the stabiliser formalism.

% ─────────────────────────────────────────────────────────────────────────────
\section*{Acknowledgements}

[To be completed.]

% ─────────────────────────────────────────────────────────────────────────────
\begin{thebibliography}{10}

\bibitem{gottesman1997stabilizer}
D.~Gottesman,
\textit{Stabilizer Codes and Quantum Error Correction},
Ph.D.\ thesis, California Institute of Technology (1997).
\href{https://arxiv.org/abs/quant-ph/9705052}{arXiv:quant-ph/9705052}

\bibitem{aaronson2004improved}
S.~Aaronson and D.~Gottesman,
\textit{Improved Simulation of Stabilizer Circuits},
Phys.\ Rev.\ A \textbf{70}, 052328 (2004).
\href{https://arxiv.org/abs/quant-ph/0406196}{arXiv:quant-ph/0406196}

\bibitem{lieb1973proof}
E.~H. Lieb and M.~B. Ruskai,
\textit{Proof of the Strong Subadditivity of Quantum-Mechanical Entropy},
J.\ Math.\ Phys.\ \textbf{14}, 1938 (1973).

\bibitem{nielsen2000quantum}
M.~A. Nielsen and I.~L. Chuang,
\textit{Quantum Computation and Quantum Information},
Cambridge University Press (2000).

\bibitem{wilde2017quantum}
M.~M. Wilde,
\textit{Quantum Information Theory}, 2nd ed.,
Cambridge University Press (2017).
\href{https://arxiv.org/abs/1106.1445}{arXiv:1106.1445}

\bibitem{sim2019expressibility}
S.~Sim, P.~D. Johnson, and A.~Aspuru-Guzik,
\textit{Expressibility and Entangling Capability of Parameterized Quantum
Circuits for Hybrid Quantum-Classical Algorithms},
Adv.\ Quantum Technol.\ \textbf{2}, 1900070 (2019).
\href{https://arxiv.org/abs/1905.10876}{arXiv:1905.10876}

\bibitem{mcclean2018barren}
J.~R. McClean, S.~Boixo, V.~N. Smelyanskiy, R.~Babbush, and H.~Neven,
\textit{Barren Plateaus in Quantum Neural Network Training Landscapes},
Nat.\ Commun.\ \textbf{9}, 4812 (2018).
\href{https://arxiv.org/abs/1803.11173}{arXiv:1803.11173}

\end{thebibliography}

\bibitem{fattal2004entanglement}
D.~Fattal, T.~S. Cubitt, Y.~Yamamoto, S.~Bravyi, and I.~L. Chuang,
\textit{Entanglement in the Stabilizer Formalism},
(2004).
\href{https://arxiv.org/abs/quant-ph/0406168}{arXiv:quant-ph/0406168}

\bibitem{audenaert2005entanglement}
K.~M.~R. Audenaert and M.~B. Plenio,
\textit{Entanglement on Mixed Stabilizer States: Normal Forms and
Reduction Procedures},
New J.\ Phys.\ \textbf{7}, 170 (2005).
\href{https://arxiv.org/abs/quant-ph/0505036}{arXiv:quant-ph/0505036}

\bibitem{bataille2022cnot}
M.~Bataille,
\textit{Quantum Circuits of CNOT Gates: Optimization and Entanglement},
Quantum Inf.\ Process.\ \textbf{21}, 269 (2022).
\href{https://doi.org/10.1007/s11128-022-03577-8}{doi:10.1007/s11128-022-03577-8}

\bibitem{goodenough2024stabilizer}
K.~Goodenough, S.~Wehner, and N.~Datta,
\textit{Bipartite Entanglement of Noisy Stabilizer States Through the Lens
of Stabilizer Codes},
IEEE Int.\ Symp.\ Inf.\ Theory (ISIT), 2024.
\href{https://arxiv.org/abs/2406.02427}{arXiv:2406.02427}

\bibitem{patel2008optimal}
K.~N. Patel, I.~L. Markov, and J.~P. Hayes,
\textit{Optimal Synthesis of Linear Reversible Circuits},
Quantum Inf.\ Comput.\ \textbf{8}(3), 282--294 (2008).

\end{thebibliography}

\end{document}
